\section{系统视角：Manager、Worker 与 Loop 的分工}
理解 veRL 的 Agentic Loop，不只是记住几个类名，而是要看清它如何把多轮推理、工具调用和服务管理拆开。Manager 负责 orchestration，Worker 负责执行，LoopBase 负责定义抽象接口，这种分层使系统可以在不同任务之间复用基础设施。

\begin{knowledgebox}{为什么这种架构适合 Agentic RL}
\begin{itemize}
  \item 多轮循环天然需要统一的状态管理
  \item 工具调用与模型推理通常要分模块复用
  \item 训练时要同时照顾 rollout、server 与任务逻辑
\end{itemize}
\end{knowledgebox}

\subsection{本章小结}
Agentic Loop 的重点不是某个类写得多复杂，而是它把 think-act-observe 这套交互逻辑抽象成了可复用的系统骨架。

\newpage
